\chapter{Modèles linéaires sous contrainte}\label{ch:qm:linear-models-under-stress}

Un modèle factoriel régresse les rendements quotidiens d'une action sur cinq facteurs, mois par mois. Deux des facteurs évoluent
ensemble, avec une corrélation de 0.97. Entre le premier et le deuxième mois, le bêta de l'action au premier facteur passe de 0.52 à $-1.59$ et son
bêta au second de 0.84 à 2.76~; leur somme ne varie que de 1.36 à 1.17, et les rendements ajustés du deuxième mois, recalculés
avec les bêtas du premier mois, conservent une corrélation de 0.92 avec leur propre ajustement. La régression n'est pas fausse~: on lui a
posé une question, celle de la répartition de l'exposition entre deux facteurs presque identiques, à laquelle trois semaines de données ne peuvent répondre. Ce chapitre porte
sur les moindres carrés lorsque leurs hypothèses sont mises à rude épreuve~: la géométrie qui les explique, la colinéarité, les pénalités qui les stabilisent,
les régresseurs mesurés avec erreur, et les \omterm{def:qm:estimation:estimator}{erreurs types} des régressions menées sur des \omterm{def:qm:linear-models-under-stress:panel}{données de panel} d'actions et de mois.

\section{La géométrie des moindres carrés}

\begin{definition}[Moindres carrés ordinaires, matrice chapeau]\label{def:qm:linear-models-under-stress:ols}
Pour $y \in \R^n$ et un plan d'expérience $X \in \R^{n \times k}$ de rang plein, les \emph{moindres carrés ordinaires}\index{moindres carrés ordinaires} (OLS) choisissent
$\hat\beta = (X^\top X)^{-1}X^\top y$, le minimiseur de $\lvert y - X\beta\rvert^2$. Les valeurs ajustées sont $\hat y = Hy$ avec la \emph{matrice
chapeau}\index{matrice chapeau} $H = X(X^\top X)^{-1}X^\top$, la projection orthogonale sur l'espace engendré par les colonnes de $X$~; ses éléments diagonaux
$h_{ii}$, dont la somme vaut $k$, sont les leviers des observations.
\end{definition}

Le résidu $y - \hat y = (I - H)y$ est orthogonal à chaque colonne de $X$. Une observation de levier proche de un tire l'ajustement
vers elle~: une mauvaise valeur dans un régresseur est le pire type de valeur aberrante, car elle est éloignée des autres dans la direction
que la régression examine.

\begin{theorem}[Gauss--Markov]\label{thm:qm:linear-models-under-stress:gm}
Si $y = X\beta + \varepsilon$ avec $\E[\varepsilon \mid X] = 0$ et $\Var(\varepsilon \mid X) = \sigma^2I$, alors $\hat\beta$ est sans biais, de
variance $\sigma^2(X^\top X)^{-1}$, et tout autre \omterm{def:qm:estimation:estimator}{estimateur} linéaire sans biais $Cy$ a une variance supérieure, d'une matrice semi-définie positive.
\end{theorem}

\begin{proof}
Écrivons $C = (X^\top X)^{-1}X^\top + D$. L'absence de biais pour tout $\beta$ impose $DX = 0$. Alors $\Var(Cy) = \sigma^2CC^\top = \sigma^2(X^\top X)^{-1} +
\sigma^2DD^\top$, les termes croisés s'annulant car $DX = 0$.
\end{proof}

Le théorème porte sur les \omterm{def:qm:estimation:estimator}{estimateurs sans biais}~; le reste du chapitre explique pourquoi un \omterm{def:qm:estimation:estimator}{estimateur} biaisé est souvent meilleur, et
ce qui se passe lorsque $\Var(\varepsilon) \ne \sigma^2I$ ou lorsque le régresseur lui-même est bruité.

\begin{theorem}[Frisch--Waugh--Lovell]\label{thm:qm:linear-models-under-stress:fwl}
Dans la régression de $y$ sur $(X_1, X_2)$, le coefficient de $X_2$ est égal au coefficient obtenu en régressant $M_1y$ sur $M_1X_2$, où $M_1 =
I - X_1(X_1^\top X_1)^{-1}X_1^\top$ retire la part expliquée par $X_1$~; les résidus des deux régressions coïncident.
\end{theorem}

\begin{proof}
À partir de $y = X_1\hat\beta_1 + X_2\hat\beta_2 + e$ avec $e$ orthogonal aux deux blocs, appliquons $M_1$~: $M_1y = M_1X_2\hat\beta_2 + e$, puisque $M_1X_1 = 0$ et
$M_1e = e$. Comme $e$ est orthogonal à $M_1X_2$, c'est la décomposition par moindres carrés de $M_1y$ sur $M_1X_2$.
\end{proof}

Un bêta «~contrôlé~» pour le marché est le bêta du rendement résiduel de l'action par rapport au marché sur le rendement résiduel du facteur par rapport au marché. Le
théorème est aussi le moyen le plus rapide de voir la colinéarité~: le coefficient de $X_2$ s'estime à partir de ce qui reste de $X_2$ après $X_1$, et
lorsque peu de chose reste, l'estimation repose sur peu de chose.

\section{Colinéarité}

\begin{definition}[Multicolinéarité, facteur d'inflation de la variance]\label{def:qm:linear-models-under-stress:vif}
La \emph{multicolinéarité}\index{multicolinéarité} est une dépendance quasi linéaire entre les colonnes du plan d'expérience. Le \emph{facteur
d'inflation de la variance}\index{facteur d'inflation de la variance} du régresseur $j$ vaut $\mathrm{VIF}_j = 1/(1 - R_j^2)$, où $R_j^2$ est le $R^2$ de la régression de $x_j$ sur les
autres régresseurs.
\end{definition}

D'après le théorème de Frisch--Waugh--Lovell, $\Var(\hat\beta_j) = \sigma^2/(n\,\widehat{\Var}(x_j)(1 - R_j^2))$~: la variance d'un coefficient
est la variance qu'il aurait seul, multipliée par son VIF. Deux facteurs de corrélation 0.97 ont $\mathrm{VIF} = 1/(1 - 0.97^2) = 16.9$ (18.5 dans
l'échantillon du chapitre), de sorte que chaque bêta est quatre fois plus bruité qu'il ne le serait seul~; la somme des deux bêtas, l'exposition
à leur direction commune, n'est pas du tout gonflée.

Sur 24 mois simulés de 21 jours, les vrais bêtas de l'action aux deux facteurs colinéaires valent 0.6 et 0.4. Les estimations mensuelles ont
des écarts types de 1.21 et 1.23 d'un mois à l'autre, un signe opposé au vrai dans 6 et 8 mois, tandis que leur somme a un écart type de 0.23
(\cref{fig:qm:linear-models-under-stress:betas})~; les bêtas aux trois facteurs non corrélés varient de 0.26 à 0.32. Un rapport de risque qui
liste les deux bêtas séparément rapporte du bruit~; celui qui rapporte leur somme, ou projette les deux facteurs sur leur première composante principale,
rapporte l'exposition.

\begin{omfigure}
\begin{tikzpicture}
\begin{axis}[width=11cm,height=5cm,xlabel={mois},ylabel={bêta OLS mensuel},
  tick label style={font=\small},label style={font=\small},xmin=0.5,xmax=24.5,ymin=-3.5,ymax=4,grid=major,grid style={gray!25},
  legend columns=3,legend style={font=\footnotesize,at={(0.5,-0.3)},anchor=north,draw=none,fill=none,
  /tikz/every even column/.append style={column sep=8pt}},table/col sep=comma]
\addplot[omDef,thick,mark=*,mark size=1.3pt] table[x=month,y=b1]{figdata/methods/16-linear-models-under-stress/betas.csv};
\addplot[omThm,thick,mark=square*,mark size=1.3pt] table[x=month,y=b2]{figdata/methods/16-linear-models-under-stress/betas.csv};
\addplot[black,very thick] table[x=month,y=sum]{figdata/methods/16-linear-models-under-stress/betas.csv};
\legend{bêta au facteur 1,bêta au facteur 2,somme des deux}
\end{axis}
\end{tikzpicture}

\omcaption{Les bêtas d'une action à deux facteurs de corrélation 0.97, estimés chaque mois à partir de 21 rendements quotidiens (vraies valeurs 0.6 et 0.4). Les
bêtas individuels oscillent entre $-2.7$ et $3.3$~; leur somme reste proche de sa vraie valeur de 1.0 (écart type 0.23). Données~: le
tutoriel du chapitre, graine fixée.}\label{fig:qm:linear-models-under-stress:betas}
\end{omfigure}

\section{Ridge, lasso et elastic net}

\begin{definition}[Régularisation, ridge, lasso, elastic net]\label{def:qm:linear-models-under-stress:penalties}
La \emph{régularisation}\index{régularisation} ajoute au critère des moindres carrés une pénalité sur la taille des coefficients. La \emph{régression
ridge}\index{régression ridge} minimise $\lvert y - X\beta\rvert^2 + \lambda\lvert\beta\rvert_2^2$~; le \emph{lasso}\index{lasso} minimise
$\frac1{2n}\lvert y - X\beta\rvert^2 + \lambda\lvert\beta\rvert_1$~; l'\emph{elastic net}\index{elastic net} utilise la pénalité $\lambda\bigl(\alpha\lvert\beta\rvert_1
+ \frac{1 - \alpha}2\lvert\beta\rvert_2^2\bigr)$, $\alpha \in [0, 1]$.
\end{definition}

\begin{proposition}[L'estimateur ridge rétrécit le long des directions singulières]\label{prop:qm:linear-models-under-stress:ridge}
Si $X = UDV^\top$ est une décomposition en valeurs singulières réduite, l'\omterm{def:qm:estimation:estimator}{estimateur} ridge vaut $\hat\beta_\lambda = \sum_jv_j\frac{d_j}{d_j^2 + \lambda}u_j^\top y$~:
le coefficient des moindres carrés selon chaque vecteur singulier droit $v_j$ est multiplié par $d_j^2/(d_j^2 + \lambda)$.
\end{proposition}

\begin{proof}
$\hat\beta_\lambda = (X^\top X + \lambda I)^{-1}X^\top y = V(D^2 + \lambda I)^{-1}DU^\top y$, et les moindres carrés correspondent au cas $\lambda = 0$, avec le facteur
$1/d_j$ selon $v_j$.
\end{proof}

Les directions colinéaires ont de petites valeurs singulières, et ce sont celles que la \omterm{def:qm:linear-models-under-stress:penalties}{régression ridge} rétrécit~: pour les deux facteurs ci-dessus, la différence $F_1 - F_2$ a
une petite valeur singulière et son coefficient est ramené vers zéro, tandis que la somme est laissée presque intacte. La \omterm{def:qm:linear-models-under-stress:penalties}{régression ridge} est la moyenne a posteriori sous une \omterm{def:qm:bayesian-methods:prior}{loi a priori} normale
sur $\beta$ (chapitre~14). La pénalité en valeur absolue du \omterm{def:qm:linear-models-under-stress:penalties}{lasso} présente un coin en zéro~: pour un plan orthonormé, sa solution est le
coefficient des moindres carrés soumis à un seuillage doux, $\operatorname{sign}(b)\max(|b| - \lambda, 0)$, de sorte qu'il annule exactement les petits coefficients et
sélectionne les variables. Au sein d'un bloc de prédicteurs corrélés, il tend à en retenir un et à écarter les autres, de façon quelque peu arbitraire~; la partie
ridge de l'\omterm{def:qm:linear-models-under-stress:penalties}{elastic net} répartit le poids sur le bloc. Le \omterm{def:qm:linear-models-under-stress:penalties}{lasso} et l'\omterm{def:qm:linear-models-under-stress:penalties}{elastic net} se calculent par descente de coordonnées cyclique, chaque coordonnée étant un
seuillage doux unidimensionnel (Friedman, Hastie et Tibshirani, 2010).

\begin{definition}[Validation croisée]\label{def:qm:linear-models-under-stress:cv}
La \emph{validation croisée}\index{validation croisée} consiste à choisir une pénalité en ajustant le modèle sur une partie des données et en mesurant l'erreur de prédiction sur le reste,
sur plusieurs découpages. Avec des séries temporelles, chaque ensemble d'apprentissage doit précéder son ensemble de test, avec un écart (une purge) au moins aussi long que la dépendance
entre les cibles, afin qu'aucune information de la période de test ne fuie dans l'ajustement.
\end{definition}

Le tutoriel prédit un rendement quotidien à partir de trente prédicteurs répartis en trois blocs de dix, corrélés à 0.8 au sein de chaque bloc~; cinq prédicteurs portent
un signal, et ce signal est faible (un modèle parfait expliquerait 5.0~\% de la variance hors échantillon). Les 1\,500 premiers jours sont découpés en
cinq plis croissants avec une purge de cinq jours, et les 500 derniers jours forment le test. Les \omterm{def:qm:linear-models-under-stress:ols}{moindres carrés ordinaires} sur l'ensemble des trente prédicteurs expliquent 1.5~\% de la
variance de test~; la \omterm{def:qm:linear-models-under-stress:penalties}{régression ridge} à la pénalité retenue par \omterm{def:qm:linear-models-under-stress:cv}{validation croisée} en explique 3.5~\%, le \omterm{def:qm:linear-models-under-stress:penalties}{lasso} 3.4~\%, l'\omterm{def:qm:linear-models-under-stress:penalties}{elastic net} 3.2~\%. Le \omterm{def:qm:linear-models-under-stress:penalties}{lasso} retient cinq prédicteurs, dont trois
sont les bons~; dans les deux blocs où il se trompe, il a retenu un voisin corrélé d'un prédicteur véritable
(\cref{fig:qm:linear-models-under-stress:path,fig:qm:linear-models-under-stress:cv}).

\begin{omfigure}
\begin{tikzpicture}
\begin{semilogxaxis}[width=11cm,height=5.4cm,xlabel={pénalité du lasso $\lambda$},ylabel={coefficient (standardisé)},x dir=reverse,
  tick label style={font=\small},label style={font=\small},xmin=1e-4,xmax=1,ymin=-0.2,ymax=0.2,grid=major,grid style={gray!25},
  table/col sep=comma]
\pgfplotsinvokeforeach{1,2,4,5,6,7,8,9,10,13,14,15,16,17,18,19,20,21,22,23,24,26,27,28,29}{
  \addplot[black!35,thin] table[x=lam,y=b#1]{figdata/methods/16-linear-models-under-stress/path.csv};
}
\pgfplotsinvokeforeach{0,3,11,12,25}{
  \addplot[omDef,very thick] table[x=lam,y=b#1]{figdata/methods/16-linear-models-under-stress/path.csv};
}
\addplot[omThm,dashed,thick,sharp plot] coordinates {(0.0398,-0.2) (0.0398,0.2)};
\end{semilogxaxis}
\end{tikzpicture}

\omcaption{Chemin du \omterm{def:qm:linear-models-under-stress:penalties}{lasso} des trente prédicteurs standardisés lorsque la pénalité décroît de 1 à $10^{-4}$~: les cinq signaux véritables en bleu, les
vingt-cinq prédicteurs de bruit en gris, et la pénalité retenue par \omterm{def:qm:linear-models-under-stress:cv}{validation croisée} 0.040 (tirets), où cinq coefficients sont non nuls. Données~: le
tutoriel du chapitre, avec graine fixée.}\label{fig:qm:linear-models-under-stress:path}
\end{omfigure}

\begin{omfigure}
\begin{tikzpicture}
\begin{semilogxaxis}[width=10cm,height=5cm,xlabel={pénalité $\lambda$},ylabel={EQM par validation croisée},x dir=reverse,
  tick label style={font=\small},label style={font=\small},xmin=1e-4,xmax=1,grid=major,grid style={gray!25},
  yticklabel style={/pgf/number format/fixed,/pgf/number format/precision=3},
  legend columns=3,legend style={font=\footnotesize,at={(0.5,-0.3)},anchor=north,draw=none,fill=none,
  /tikz/every even column/.append style={column sep=8pt}},table/col sep=comma]
\addplot[omDef,very thick] table[x=lam,y=ridge]{figdata/methods/16-linear-models-under-stress/cv.csv};
\addplot[omThm,very thick,dashed] table[x=lam,y=lasso]{figdata/methods/16-linear-models-under-stress/cv.csv};
\addplot[black,thick,dotted] table[x=lam,y=enet]{figdata/methods/16-linear-models-under-stress/cv.csv};
\legend{régression ridge,lasso,{elastic net, $\alpha = 0.5$}}
\end{semilogxaxis}
\end{tikzpicture}

\omcaption{Erreur de \omterm{def:qm:linear-models-under-stress:cv}{validation croisée} ordonnée dans le temps des trois régressions pénalisées en fonction de la pénalité (cinq plis croissants avec une
purge de cinq jours, sur les 1\,500 premiers jours). Chaque courbe présente un minimum intérieur~: une pénalité trop faible ajuste le bruit, une pénalité trop forte écarte le signal. Données~: le
tutoriel du chapitre, avec graine fixée.}\label{fig:qm:linear-models-under-stress:cv}
\end{omfigure}

\section{Régression pondérée et erreurs sur les variables}

\begin{definition}[Moindres carrés pondérés et généralisés]\label{def:qm:linear-models-under-stress:gls}
La méthode des \emph{moindres carrés pondérés}\index{moindres carrés pondérés} minimise $\sum_iw_i(y_i - x_i^\top\beta)^2$, avec des poids inversement proportionnels à la
variance de l'erreur de chaque observation. La méthode des \emph{moindres carrés généralisés}\index{moindres carrés généralisés} minimise $(y - X\beta)^\top\Omega^{-1}(y -
X\beta)$ pour une matrice de covariance des erreurs $\Omega$~; ce sont les \omterm{def:qm:linear-models-under-stress:ols}{moindres carrés ordinaires} appliqués après multiplication du modèle par $\Omega^{-1/2}$, et c'est l'\omterm{def:qm:estimation:estimator}{estimateur} linéaire sans biais
le plus efficace lorsque la matrice $\Omega$ est connue.
\end{definition}

Lorsque la matrice $\Omega$ est inconnue, le choix usuel en finance est celui des \omterm{def:qm:linear-models-under-stress:ols}{moindres carrés ordinaires} avec une \omterm{def:qm:estimation:estimator}{erreur type} qui ne néglige pas cette structure~: la \omterm{def:qm:estimation:sandwich}{variance sandwich}
du chapitre~11, HAC pour les séries temporelles, \omterm{def:qm:linear-models-under-stress:fm}{erreurs types groupées} (clustered) pour les \omterm{def:qm:linear-models-under-stress:panel}{données de panel} (ci-dessous).

\begin{definition}[Erreurs sur les variables, biais d'atténuation, moindres carrés totaux]\label{def:qm:linear-models-under-stress:eiv}
Un modèle à \emph{erreurs sur les variables}\index{erreurs sur les variables} observe le régresseur avec du bruit, $\tilde x = x + u$, avec $u$ indépendant de $x$
et de l'erreur de l'équation. Le rétrécissement de la pente des \omterm{def:qm:linear-models-under-stress:ols}{moindres carrés ordinaires} vers zéro qui en résulte est le \emph{biais d'atténuation}\index{biais d'atténuation}. \emph{Moindres
carrés totaux}\index{moindres carrés totaux} ajuste la droite qui minimise la somme des carrés des distances perpendiculaires, ce qui est convergent lorsque les erreurs sur
$x$ et $y$ ont des variances égales.
\end{definition}

\begin{proposition}[Le facteur d'atténuation]\label{prop:qm:linear-models-under-stress:attenuation}
Si $y = \beta x + \varepsilon$ et $\tilde x = x + u$ avec $x$, $u$, $\varepsilon$ indépendants, la pente des \omterm{def:qm:linear-models-under-stress:ols}{moindres carrés ordinaires} de $y$ sur $\tilde x$ converge vers
$\lambda\beta$ avec $\lambda = \Var(x)/(\Var(x) + \Var(u))$.
\end{proposition}

\begin{proof}
$\Cov(y, \tilde x) = \beta\Var(x)$ et $\Var(\tilde x) = \Var(x) + \Var(u)$~; la pente en est le rapport.
\end{proof}

Un desk couvre une obligation avec un future sur obligations, et estime le ratio de couverture en régressant les variations quotidiennes de prix de l'obligation sur celles du future. L'obligation
est valorisée au cours médian synchrone à la clôture~; le prix quotidien enregistré du future est sa dernière transaction avant la clôture, décalé de quelques minutes et
au cours acheteur ou au cours vendeur. On modélise le niveau enregistré comme le niveau vrai plus une erreur indépendante d'écart type 0.14~point~; la variation quotidienne vraie du
future a un écart type de 0.40, et l'obligation varie de 0.85 fois le future plus un terme idiosyncratique de 0.05. La variation enregistrée porte
la différence de deux erreurs de niveau, de variance $2 \times 0.14^2 = 0.039$, donc $\lambda = 0.16/(0.16 + 0.039) = 0.80$. Sur deux ans de données
à graine fixée, le ratio de couverture par \omterm{def:qm:linear-models-under-stress:ols}{moindres carrés ordinaires} vaut 0.68, et sur 2\,000 historiques simulés il vaut en moyenne 0.683 avec un écart type de 0.017~: précisément faux.

Deux corrections découlent de la structure de l'erreur. Une erreur de niveau rend négativement corrélées les variations enregistrées consécutives, avec une première
autocovariance $-\sigma_e^2$ (l'argument de Roll de 1984 sur la fourchette bid-ask), donc $\hat\lambda = (\hat\gamma(0) + 2\hat\gamma(1))/\hat\gamma(0)$
estime l'atténuation à partir de la série propre du future~; ici l'autocorrélation d'ordre un vaut $-0.098$, $\hat\lambda = 0.80$, et le
ratio corrigé vaut 0.85. Sur l'ensemble des historiques, cette correction est sans biais (0.862 en moyenne) mais bruitée (écart type 0.095). Autrement,
les variations hebdomadaires ajoutent cinq jours de mouvement vrai aux deux mêmes erreurs de niveau, si bien que leur atténuation n'est que de $5 \times 0.16/(5 \times 0.16 + 0.039) =
0.95$~: la régression hebdomadaire donne 0.80 sur cet historique, 0.810 en moyenne avec un écart type de 0.022
(\cref{fig:qm:linear-models-under-stress:hedge}). Les \omterm{def:qm:linear-models-under-stress:eiv}{moindres carrés totaux} donnent 0.74~: leur hypothèse de variances d'erreur égales est fausse ici.

\begin{omfigure}
\begin{tikzpicture}
\begin{axis}[width=10cm,height=5cm,xlabel={ratio de couverture estimé},ylabel={densité sur 2\,000 historiques},
  tick label style={font=\small},label style={font=\small},xmin=0.5,xmax=1.2,ymin=0,ymax=26,grid=major,grid style={gray!25},
  legend columns=4,legend style={font=\footnotesize,at={(0.5,-0.3)},anchor=north,draw=none,fill=none,
  /tikz/every even column/.append style={column sep=6pt}},table/col sep=comma]
\addplot[omThm,thick,const plot mark mid,fill=omThm!15] table[x=h,y=ols]{figdata/methods/16-linear-models-under-stress/hedge.csv};
\addplot[black,thick,const plot mark mid] table[x=h,y=roll]{figdata/methods/16-linear-models-under-stress/hedge.csv};
\addplot[omDef,very thick,const plot mark mid] table[x=h,y=weekly]{figdata/methods/16-linear-models-under-stress/hedge.csv};
\addplot[black,dashed,thin,sharp plot] coordinates {(0.85,0) (0.85,26)};
\legend{MCO quotidiens,corrigé de Roll,MCO hebdomadaires,vrai 0.85}
\end{axis}
\end{tikzpicture}

\omcaption{Distributions d'échantillonnage de trois estimations du ratio de couverture sur 2\,000 historiques simulés de deux ans~: MCO quotidiens sur les prix bruités du future
(moyenne 0.683, atténués d'un facteur 0.80), les mêmes divisés par l'atténuation estimée à partir de l'autocorrélation d'ordre un du future (moyenne 0.862, écart type
0.095), et MCO sur les variations hebdomadaires (moyenne 0.810, écart type 0.022). Données~: le tutoriel du chapitre, avec graine fixée.}\label{fig:qm:linear-models-under-stress:hedge}
\end{omfigure}

Le coût se trouve dans le résidu. Mesurée par rapport aux variations synchrones vraies, la couverture au vrai ratio laisse un résidu quotidien de 0.049
points~; au ratio MCO quotidien 0.086, soit 74~\% de plus~; au ratio corrigé 0.049~; au ratio hebdomadaire 0.053. Sans couverture, l'obligation bouge de 0.34 par jour.

\section{Régressions en coupe transversale dans le temps, et panels}

\begin{definition}[Données de panel, effets fixes]\label{def:qm:linear-models-under-stress:panel}
Les \emph{données de panel}\index{données de panel} observent de nombreuses unités (actions, fonds, contreparties) sur de nombreuses périodes. Une régression à \emph{effets fixes}\index{effets fixes}
attribue à chaque unité (ou à chaque période) sa propre ordonnée à l'origine, ce qui, par le théorème de Frisch--Waugh--Lovell, revient à centrer chaque variable
au sein de l'unité.
\end{definition}

\begin{definition}[Régression de Fama--MacBeth, erreurs types groupées]\label{def:qm:linear-models-under-stress:fm}
Une \emph{régression de Fama--MacBeth}\index{régression de Fama--MacBeth} (Fama et MacBeth, 1973) effectue une régression en coupe transversale par période et
rapporte la moyenne temporelle des pentes, avec pour \omterm{def:qm:estimation:estimator}{erreur type} leur écart type sur $\sqrt T$. Les \emph{erreurs types groupées}\index{erreurs types groupées} (clustered
standard errors) calculent la \omterm{def:qm:estimation:sandwich}{variance sandwich} en sommant les scores à l'intérieur de groupes (entreprises, ou périodes), ce qui autorise toute
corrélation à l'intérieur d'un groupe et aucune entre groupes.
\end{definition}

L'\omterm{def:qm:estimation:estimator}{erreur type} de Fama--MacBeth est correcte lorsque les pentes sont indépendantes dans le temps~; un effet de période (les résidus de toutes les actions évoluant
ensemble au cours d'un mois) est pris en compte, puisque la pente de chaque mois l'absorbe. Un effet d'entreprise persistant ne l'est pas~: si une caractéristique et le résidu
comportent chacun une composante fixe pour chaque entreprise, la pente de chaque mois partage la même erreur, que l'écart type temporel ne peut pas voir, et qu'une correction
de Newey--West à quelques retards ne voit pas davantage (Petersen, 2009). Dans un panel simulé de 100 entreprises sur 60 mois, où les effets d'entreprise représentent la moitié
de la variance de la caractéristique et du résidu, l'écart type vrai de la pente sur 500 panels vaut 0.105. Les \omterm{def:qm:estimation:estimator}{erreurs types} classiques et de White valent en moyenne 0.026, celle de Fama--MacBeth 0.023 (0.022 avec trois retards de Newey--West),
et seule l'\omterm{def:qm:estimation:estimator}{erreur type} groupée par entreprise, 0.101,
est correcte (\cref{fig:qm:linear-models-under-stress:panel}).

\begin{omfigure}
\begin{tikzpicture}
\begin{axis}[width=10cm,height=4.8cm,xbar,bar width=8pt,
  xlabel={erreur type de la pente (moyenne sur 500 panels)},
  ytick={0,1,2,3,4,5},yticklabels={écart type vrai,MCO classiques,White,groupées par entreprise,Fama--MacBeth,{Fama--MacBeth, 3 retards}},
  tick label style={font=\small},label style={font=\small},
  xmin=0,xmax=0.13,ymin=-0.6,ymax=5.6,grid=major,grid style={gray!25},
  xticklabel style={/pgf/number format/fixed},scaled x ticks=false,table/col sep=comma]
\addplot[fill=omDef!60,draw=omDef,nodes near coords,nodes near coords style={font=\footnotesize,/pgf/number format/fixed,/pgf/number format/precision=3}]
  table[x=se,y=k]{figdata/methods/16-linear-models-under-stress/panel.csv};
\end{axis}
\end{tikzpicture}

\omcaption{\omterm{def:qm:estimation:estimator}{Erreurs types} de la pente d'un panel lorsque la caractéristique et le résidu comportent tous deux une composante persistante par entreprise~: seul le regroupement
par entreprise retrouve l'écart type d'échantillonnage vrai (0.105)~; les \omterm{def:qm:estimation:estimator}{erreurs types} classiques, de White et de Fama--MacBeth le sous-estiment environ quatre fois.
Données~: le tutoriel du chapitre, avec graine fixée.}\label{fig:qm:linear-models-under-stress:panel}
\end{omfigure}

\section{Tutoriel~: le ratio de couverture qui a rétréci}\label{tut:qm:linear-models-under-stress:hedge}

\begin{tutorial}
\textbf{Objectif.} Stabiliser une régression colinéaire, sélectionner des prédicteurs par \omterm{def:qm:linear-models-under-stress:cv}{validation croisée} ordonnée dans le temps, et corriger un ratio de couverture atténué.
\textbf{État final~:} \cref{fig:qm:linear-models-under-stress:path,fig:qm:linear-models-under-stress:cv,fig:qm:linear-models-under-stress:hedge}, le
$R^2$ de test de 1.5~\% (MCO) et 3.5~\% (ridge), et les ratios de couverture 0.68, 0.85 et 0.80.
\begin{tutsteps}
\item \textbf{Ridge par la SVD et \omterm{def:qm:linear-models-under-stress:penalties}{elastic net} par descente de coordonnées.}
\omcode{code/firm/linreg/firm_linreg.py}{76}{104}{Ridge et elastic net.}\label{lst:qm:linear-models-under-stress:enet}
\item \textbf{Plis ordonnés dans le temps avec purge.}
\omcode{code/firm/linreg/firm_linreg.py}{107}{117}{Plis de validation croisée à fenêtre croissante.}\label{lst:qm:linear-models-under-stress:folds}
\item \textbf{Exécuter} \texttt{fit\_all()}, \texttt{hedge()}, \texttt{hedge\_mc()} et \texttt{panel\_truth()} dans \texttt{qm\_linreg.py}, puis
\texttt{fig\_linreg.py}.
\end{tutsteps}
\textbf{Pour aller plus loin.} Remplacer les deux facteurs colinéaires par leur somme et leur différence, et observer le bêta de la différence porter seul le
bruit~; réduire la purge à zéro avec des cibles de cinq jours qui se chevauchent, et mesurer à quel point la \omterm{def:qm:linear-models-under-stress:cv}{validation croisée} flatte l'ajustement.
\end{tutorial}

\section{Construction~: le module de régression}\label{bld:qm:linear-models-under-stress:linreg}

\begin{build}
\bfield{Rôle} Toute régression que la petite entreprise exécute (bêtas, ratios de couverture, ajustements de signaux, tests de panel) passe par un seul module dont
les \omterm{def:qm:estimation:estimator}{erreurs types} sont choisies, et non laissées par défaut.
\bfield{Interface} \texttt{ols(X, y, cov, lags, groups)} avec \texttt{cov} parmi classical, \texttt{hc}, \texttt{hac}, \texttt{cluster}~;
\texttt{vif(X)}~; \texttt{ridge(X, y, lam)}~; \texttt{elastic\_net(X, y, lam, alpha)}~; \texttt{time\_series\_folds(n, n\_folds, purge)}~;
\texttt{fama\_macbeth(ys, Xs, lags)}~; \texttt{tls(x, y)}.
\bfield{Règles} Le type de covariance est un choix obligatoire dans les rapports~; les prédicteurs sont standardisés sur le seul pli d'entraînement~; les plis n'entraînent jamais
sur des données postérieures à leur bloc de test~; réchauffement à chaud le long d'un chemin de pénalité.
\bfield{Tests d'acceptation} \texttt{code/firm/linreg/tests/}~: les \omterm{def:qm:estimation:estimator}{erreurs types} de White égalent la dispersion d'échantillonnage simulée en présence d'hétéroscédasticité~;
l'identité de Frisch--Waugh--Lovell~; la VIF de deux variables de corrélation 0.97~; ridge égale ses équations normales et l'\omterm{def:qm:linear-models-under-stress:penalties}{elastic net} avec
$\alpha = 0$ égale ridge~; le \omterm{def:qm:linear-models-under-stress:penalties}{lasso} applique un seuillage doux sur un plan d'expérience orthonormé~; plis purgés~; Fama--MacBeth sur des périodes indépendantes.
\bfield{Pour aller plus loin} Régression de Deming avec un rapport de variances d'erreur connu~; le \omterm{def:qm:linear-models-under-stress:penalties}{lasso} adaptatif~; regroupement à deux voies par entreprise et par mois.
\end{build}

\begin{omsources}
\item R.~Frisch et F.~V.~Waugh, « Partial time regressions as compared with individual trends », \emph{Econometrica} 1, 1933~; M.~C.~Lovell,
\emph{Journal of the American Statistical Association} 58, 1963.
\item A.~E.~Hoerl et R.~W.~Kennard, « Ridge regression: biased estimation for nonorthogonal problems », \emph{Technometrics} 12, 1970.
\item R.~Tibshirani, « Regression shrinkage and selection via the lasso », \emph{Journal of the Royal Statistical Society B} 58, 1996.
\item H.~Zou et T.~Hastie, « Regularization and variable selection via the elastic net », \emph{Journal of the Royal Statistical Society B}
67, 2005.
\item J.~Friedman, T.~Hastie et R.~Tibshirani, « Regularization paths for generalized linear models via coordinate descent », \emph{Journal of
Statistical Software} 33, 2010.
\item E.~F.~Fama et J.~D.~MacBeth, « Risk, return, and equilibrium: empirical tests », \emph{Journal of Political Economy} 81, 1973.
\item M.~A.~Petersen, « Estimating standard errors in finance panel data sets: comparing approaches », \emph{Review of Financial Studies} 22, 2009.
\item R.~Roll, « A simple implicit measure of the effective bid-ask spread in an efficient market », \emph{Journal of Finance} 39, 1984.
\end{omsources}

\section{Exercices}

\begin{exercise}[$\star$]\label{exo:qm:linear-models-under-stress:1}
Deux régresseurs ont une corrélation de 0.9. Quel est le \omterm{def:qm:linear-models-under-stress:vif}{facteur d'inflation de la variance}, et de quel facteur l'\omterm{def:qm:estimation:estimator}{erreur type} de chaque coefficient augmente-t-elle~?
\end{exercise}

\begin{exercise}[$\star$]\label{exo:qm:linear-models-under-stress:2}
Un régresseur a une variance vraie de 1 et il est mesuré avec un bruit indépendant de variance 0.25. Si la pente vraie vaut 2, qu'estime la MCO~?
\end{exercise}

\begin{exercise}[$\star$]\label{exo:qm:linear-models-under-stress:3}
Montrer que les leviers $h_{ii}$ somment à $k$, le nombre de colonnes de $X$.
\end{exercise}

\begin{exercise}[$\star\star$]\label{exo:qm:linear-models-under-stress:4}
Pour un plan orthonormé ($X^\top X = nI$), dériver la solution du \omterm{def:qm:linear-models-under-stress:penalties}{lasso} coordonnée par coordonnée.
\end{exercise}

\begin{exercise}[$\star\star$]\label{exo:qm:linear-models-under-stress:5}
Les variations quotidiennes enregistrées d'un contrat à terme ont une variance de 0.199 et une première autocovariance $-0.0196$. Estimer la variance de l'erreur de niveau et le
facteur d'atténuation d'une régression de couverture quotidienne, ainsi que celui d'une régression hebdomadaire.
\end{exercise}

\begin{exercise}[$\star\star$]\label{exo:qm:linear-models-under-stress:6}
Avec les deux facteurs colinéaires du chapitre ($\rho = 0.97$, variances unitaires), quelle combinaison linéaire des deux bêtas les données déterminent-elles,
et de combien est-elle plus précise que chacun des bêtas~?
\end{exercise}

\begin{exercise}[$\star\star\star$]\label{exo:qm:linear-models-under-stress:7}
\emph{Programmation.} Simuler le panel du chapitre avec effets par entreprise et comparer la moyenne des \omterm{def:qm:estimation:estimator}{erreurs types} classiques, groupées (clustered) et de Fama--MacBeth avec
l'écart-type d'échantillonnage vrai de la pente sur 500 panels.
\end{exercise}

\begin{exercise}[$\star\star\star$]\label{exo:qm:linear-models-under-stress:8}
\emph{Repérer le défaut.} « Nous avons choisi la pénalité du \omterm{def:qm:linear-models-under-stress:penalties}{lasso} par \omterm{def:qm:linear-models-under-stress:cv}{validation croisée} à dix plis aléatoires sur cinq ans de données quotidiennes, en prédisant pour chaque
action son rendement à vingt jours~; le $R^2$ validé est de 8~\%. »
\end{exercise}

\section{Problème~: la couverture qui s'est réduite}

\begin{problem}[{Problème du week-end --- couverture d'une obligation par un contrat à terme aux prix bruités}]\label{pb:qm:linear-models-under-stress:1}
La variation quotidienne du prix d'une obligation vaut 0.85 fois la variation vraie du contrat à terme plus un écart idiosyncratique de 0.05 point~; la variation vraie du contrat à terme a un écart-type
de 0.40 point~; son prix quotidien enregistré porte une erreur de niveau indépendante d'écart-type 0.14. Deux ans (500 jours) des
données simulées du chapitre sont disponibles.

\medskip\noindent\textbf{Partie I --- La couverture naïve.}
\begin{enumerate}
\item Quel ratio de couverture donne la MCO quotidienne~?
\item Quel facteur d'atténuation le modèle prédit-il, et qu'implique-t-il pour le ratio MCO moyen~?
\item Sur 2\,000 historiques simulés, quels sont la moyenne et l'écart-type du ratio MCO quotidien~?
\item Pourquoi une estimation précise n'est-elle pas correcte ici~?
\item Quel risque quotidien résiduel le ratio MCO laisse-t-il, par rapport à celui du ratio vrai~?
\end{enumerate}

\medskip\noindent\textbf{Partie II --- Corrections.}
\begin{enumerate}[resume]
\item Quelle est la première autocorrélation des variations enregistrées du contrat à terme, et qu'implique-t-elle pour $\hat\lambda$~?
\item Quel est le ratio corrigé, et quel est son risque résiduel~?
\item Quelle atténuation subit une régression hebdomadaire, et quel ratio donne-t-elle~?
\item Quelles sont la moyenne et l'écart-type des ratios corrigé et hebdomadaire sur les historiques~?
\item Que donnent les \omterm{def:qm:linear-models-under-stress:eiv}{moindres carrés totaux}, et pourquoi sont-ils faux ici~?
\end{enumerate}

\medskip\noindent\textbf{Partie III --- Au-delà de la couverture.}
\begin{enumerate}[resume]
\item Quelle serait l'atténuation avec une erreur de niveau de 0.07 au lieu de 0.14~?
\item Comment des prix synchrones (relevés au même instant, au milieu de la fourchette) modifieraient-ils le problème~?
\item Si le ratio de couverture était estimé sur des variations mensuelles, que gagnerait-on et que perdrait-on~?
\item La colinéarité du modèle factoriel du chapitre pourrait-elle être réparée par les mêmes moyens~?
\item Comment vérifieriez-vous le ratio estimé avant de le trader~?
\end{enumerate}

\medskip\noindent\textbf{Partie IV --- Jugement.}
\begin{enumerate}[resume]
\item Quel ratio utiliseriez-vous, et pourquoi~?
\item Que devrait changer l'équipe de données~?
\item Quand l'atténuation de la pente d'une régression est-elle sans conséquence~?
\item Énoncer le \emph{résultat nommé}~: le facteur d'atténuation, le ratio corrigé et le risque résiduel que laisse chacun.
\item En une phrase~: qu'est-ce que le bruit dans un régresseur fait que le bruit dans la variable dépendante ne fait pas~?
\end{enumerate}
\end{problem}

\section{Questions d'entretien}

\begin{interviewq}[$\star$ \iqroles{chercheur, ingénieur ML}]\label{iq:qm:linear-models-under-stress:1}
Deux de vos régresseurs ont une corrélation de 0.97. Que deviennent leurs coefficients, et que faites-vous~?
\end{interviewq}

\begin{interviewq}[$\star\star$ \iqroles{chercheur, trader}]\label{iq:qm:linear-models-under-stress:2}
Votre ratio de couverture par régression vaut 0.68, mais le ratio de duration indique 0.85. Qu'est-ce qui pourrait expliquer l'écart~?
\end{interviewq}

\begin{interviewq}[$\star\star$ \iqroles{ingénieur ML, chercheur}]\label{iq:qm:linear-models-under-stress:3}
Ridge ou \omterm{def:qm:linear-models-under-stress:penalties}{lasso}~: que fait chacun aux prédicteurs corrélés~?
\end{interviewq}

\begin{interviewq}[$\star\star$ \iqroles{ingénieur ML, chercheur}]\label{iq:qm:linear-models-under-stress:4}
Comment effectuez-vous une \omterm{def:qm:linear-models-under-stress:cv}{validation croisée} d'un modèle qui prédit des rendements à vingt jours à partir de données quotidiennes~?
\end{interviewq}

\begin{interviewq}[$\star\star$ \iqroles{chercheur}]\label{iq:qm:linear-models-under-stress:5}
Quand les \omterm{def:qm:estimation:estimator}{erreurs types} de Fama--MacBeth sont-elles fausses, et lesquelles utilisez-vous à la place~?
\end{interviewq}

\begin{interviewq}[$\star\star\star$ \iqroles{chercheur, ingénieur ML}]\label{iq:qm:linear-models-under-stress:6}
Énoncer et démontrer le théorème de Frisch--Waugh--Lovell, et en donner une utilisation sur un desk de trading.
\end{interviewq}
